
\documentclass[12pt,a4paper]{article}

% Formato general APA 7 (adaptado a manuscrito académico con límite de páginas)
\usepackage[spanish,es-tabla]{babel}
\usepackage[utf8]{inputenc}
\usepackage[T1]{fontenc}
\usepackage{mathptmx}
\usepackage[a4paper,margin=2.54cm]{geometry}
\usepackage{setspace}
\setstretch{1.0}
\usepackage{graphicx}
\usepackage{tabularx}
\usepackage{array}
\usepackage{booktabs}
\usepackage{enumitem}
\usepackage{float}
\usepackage{caption}
\usepackage{xurl}
\usepackage{hyperref}
\usepackage{xcolor}
\usepackage{titlesec}
\usepackage{microtype}
\usepackage{longtable}
\usepackage{fancyhdr}
\usepackage{etoolbox}

\hypersetup{colorlinks=true, linkcolor=black, citecolor=black, urlcolor=blue}
\sloppy

% Encabezado y paginación APA: número de página arriba a la derecha.
\pagestyle{fancy}
\fancyhf{}
\rhead{\thepage}
\renewcommand{\headrulewidth}{0pt}

% Sangría APA y espaciado de párrafos.
\setlength{\parindent}{1.27cm}
\setlength{\parskip}{0pt}
\setlist[itemize]{leftmargin=1.27cm,itemsep=0.15em,topsep=0.2em}
\setlist[enumerate]{leftmargin=1.27cm,itemsep=0.15em,topsep=0.2em}

% Encabezados estilo APA simplificado.
\titleformat{\section}{\centering\bfseries\normalsize}{\thesection.}{0.5em}{}
\titleformat{\subsection}{\bfseries\normalsize}{\thesubsection.}{0.5em}{}
\titleformat{\subsubsection}{\bfseries\itshape\normalsize}{\thesubsubsection.}{0.5em}{}

% Tablas y figuras con formato APA: número y título encima.
\captionsetup[table]{labelfont=bf,textfont=it,labelsep=newline,justification=centering,singlelinecheck=false}
\captionsetup[figure]{labelfont=bf,textfont=it,labelsep=newline,justification=centering,singlelinecheck=false}
\renewcommand{\tablename}{Tabla}
\renewcommand{\figurename}{Figura}
\renewcommand{\arraystretch}{1.12}
\newcolumntype{Y}{>{\raggedright\arraybackslash}X}

% Referencias APA manuales con sangría francesa.
\newcommand{\apaRef}[1]{\par\noindent\hangindent=1.27cm\hangafter=1 #1\par}

\begin{document}

\begin{titlepage}
\thispagestyle{fancy}
\centering
\vspace*{3cm}
{\bfseries\Large Predicción y modelado del burnout académico en estudiantes universitarios mediante técnicas de Ciencia de Datos: una revisión sistemática de literatura\par}
\vspace{2cm}
Felipe Martínez González\par
Ciencia de Datos\par
Metodología de la Investigación en Ciencia de Datos\par
Universidad Tecnológica Metropolitana\par
fmartinezgo@utem.cl\par
\vspace{2cm}
06 de mayo de 2026\par
\end{titlepage}

\section*{Resumen}
\addcontentsline{toc}{section}{Resumen}
El burnout académico constituye un problema relevante en educación superior, debido a su relación con el agotamiento emocional, el desenganche, el bajo bienestar psicológico, el rendimiento académico y la permanencia estudiantil. Desde la Ciencia de Datos, este fenómeno puede abordarse mediante modelos capaces de identificar factores de riesgo, patrones de comportamiento y perfiles de estudiantes que podrían requerir apoyo temprano. El objetivo de esta revisión sistemática fue analizar la literatura científica sobre el uso de técnicas de Ciencia de Datos para la predicción o modelado del burnout académico en estudiantes universitarios. Se realizó una búsqueda en Web of Science y Scopus, seguida de deduplicación, cribado, recuperación de texto completo y selección final mediante PRISMA. Se identificaron 1.020 registros, de los cuales 2 fueron eliminados por duplicidad. Se cribaron 1.018 registros, se excluyeron 1.004, se buscaron 14 informes para recuperación, 2 no fueron recuperados, 12 fueron evaluados para elegibilidad y finalmente 6 estudios fueron incluidos. Los métodos identificados incluyeron regresión lineal múltiple, análisis multinivel, modelos de ecuaciones estructurales, análisis de redes, minería de datos institucional y clustering. Se concluye que la Ciencia de Datos presenta alto potencial para apoyar la detección temprana del burnout académico; sin embargo, el campo requiere mayor reproducibilidad, validación externa, integración de datos institucionales y desarrollo de modelos predictivos explicables, éticos y aplicables en educación superior.

\noindent\textit{Palabras clave:} burnout académico, educación superior, estrés académico, detección temprana, modelado predictivo.

\section{Introducción}

El burnout académico es un problema relevante en educación superior por su relación con agotamiento emocional, desenganche, bajo bienestar, rendimiento académico y permanencia estudiantil (Cuevas Caravaca et al., 2025; Mud\l{}o-G\l{}agolska \& Larionow, 2025). Este fenómeno se asocia con exigencias académicas sostenidas, presión por el rendimiento, adaptación universitaria, dificultades económicas y falta de mecanismos institucionales de detección temprana.

Desde la Ciencia de Datos, el burnout académico puede abordarse mediante modelos capaces de identificar patrones de riesgo a partir de variables psicológicas, académicas, sociodemográficas y conductuales. Técnicas como regresión, clasificación, clustering, análisis de redes, modelos longitudinales o minería de datos institucional pueden apoyar la detección temprana y orientar intervenciones preventivas (Alluri et al., 2026; Bojorque et al., 2025).

Sin embargo, la literatura presenta heterogeneidad en instrumentos de medición, tipos de datos, métodos analíticos y métricas de evaluación. Además, no siempre existe disponibilidad de datos, código o validación externa, lo que limita la reproducibilidad y transferencia de modelos a otros contextos universitarios.

El objetivo de esta revisión es analizar sistemáticamente la literatura científica sobre el uso de técnicas de Ciencia de Datos para la predicción o modelado del burnout académico en estudiantes universitarios.

La pregunta principal de investigación es: ¿qué técnicas de Ciencia de Datos se han utilizado para predecir o modelar el burnout académico en estudiantes universitarios y qué vacíos metodológicos permanecen en la literatura? De forma específica, se busca identificar los tipos de datos utilizados, los métodos computacionales aplicados, las métricas reportadas y las limitaciones que afectan la reproducibilidad y aplicabilidad de los modelos.

\section{Metodología}

Se realizó una revisión sistemática de literatura orientada a identificar estudios sobre predicción o modelado del burnout académico en estudiantes universitarios mediante técnicas de Ciencia de Datos. El proceso consideró identificación, deduplicación, cribado, recuperación de texto completo, evaluación de elegibilidad e inclusión final, siguiendo el flujo PRISMA 2020 (Page et al., 2021).

\subsection{Estrategia de búsqueda}

\begin{itemize}

    \item \textbf{Bases de datos consultadas:} Web of Science (WoS) y Scopus.

    \item \textbf{Fecha de búsqueda:} 24 de abril de 2026.

    \item \textbf{Palabras clave y sinónimos:}

    \begin{itemize}
        \item \textit{Problema:} Academic Burnout, Student Burnout, Burnout, Emotional Exhaustion.
        
        \item \textit{Palabras relacionadas:} University Students, College Students, Undergraduate, Higher Education.
        
        \item \textit{Metodos:} Machine Learning, Data Science, Prediction, Predictive Model, Classification, Regression, Clustering.
    \end{itemize}

    \item \textbf{Operadores booleanos:} Se utilizaron AND para conectar las diferentes dimensiones temáticas y OR para agrupar los sinónimos dentro de cada dimensión.

    \item \textbf{Cadenas de búsqueda completas:}
    \begin{quote}\
    (\texttt{Academic burnout OR Student mental health OR Higher education}) AND 
    (\texttt{Academic stress OR Risk factors}) AND 
    (\texttt{Machine learning OR Predictive modeling OR Early detection})
    \end{quote}

    \item \textbf{Ajustes por base de datos:} Tanto en WoS como en Scopus se filtraron los papers desde el año 2021 hasta el 2026, pero solo en WoS se seleccionaron papers en español para reducir la cantidad.

\end{itemize}

\subsection{Criterios de inclusión y exclusión}

Se aplicaron criterios rigurosos para filtrar los registros iniciales y llegar a la selección
final de los 6 artículos analizados en esta revisión. Estos criterios se detallan en la Tabla~\ref{tab:criterios_inclusion_exclusion}.

\begin{table}[htbp]
\centering
\caption{Criterios de inclusión y exclusión}
\label{tab:criterios_inclusion_exclusion}
\renewcommand{\arraystretch}{1.15}
\small

\begin{tabularx}{\textwidth}{|
>{\raggedright\arraybackslash}p{2.3cm}|
>{\raggedright\arraybackslash}X|
>{\raggedright\arraybackslash}X|}
\hline
\textbf{Tipo} & \textbf{Criterio} & \textbf{Justificación} \\
\hline

Inclusión &
Artículos revisados por pares. &
Asegura calidad científica mínima. \\
\hline

Inclusión &
Estudios sobre burnout académico, agotamiento académico o burnout estudiantil. &
Asegura pertinencia temática. \\
\hline

Inclusión &
Estudios en estudiantes universitarios o educación superior. &
Mantiene coherencia con la población objetivo. \\
\hline

Inclusión &
Estudios con métodos cuantitativos, computacionales o de Ciencia de Datos. &
Permite analizar técnicas de modelado o predicción. \\
\hline

Inclusión &
Estudios con resumen, metadatos o texto completo disponible. &
Permite evaluar pertinencia metodológica. \\
\hline

Exclusión &
Registros duplicados. &
Evita sobrecontar evidencia. \\
\hline

Exclusión &
Estudios sin relación directa con burnout académico. &
Reduce ruido temático. \\
\hline

Exclusión &
Estudios solo sobre salud mental general sin vínculo con burnout. &
Evita desviar el alcance de la revisión. \\
\hline

Exclusión &
Población no universitaria sin aporte metodológico relevante. &
Mantiene coherencia poblacional. \\
\hline

Exclusión &
Documentos sin análisis estadístico, computacional o predictivo. &
No aportan al enfoque de Ciencia de Datos. \\
\hline

Exclusión &
Texto completo no recuperado. &
No permite evaluación adecuada. \\
\hline

\end{tabularx}
\end{table}

\subsection{Selección de estudios mediante PRISMA}

El proceso PRISMA permitió pasar de 1.020 registros iniciales a 6 estudios incluidos. Luego de eliminar 2 duplicados, se cribaron 1.018 registros. En el screening se excluyeron 1.004 por no cumplir criterios temáticos, poblacionales o metodológicos. Se buscaron 14 informes para recuperación; 2 no fueron recuperados. Se evaluaron 12 informes a texto completo y se excluyeron 6 con razón, quedando 6 estudios en la síntesis final.

\begin{figure}[H]
\centering
\includegraphics[width=0.78\linewidth]{figuras/prisma_burnout_final.jpg}
\caption{Diagrama de flujo PRISMA 2020 del proceso de selección de estudios.}
\label{fig:prisma}
\end{figure}

\subsection{Extracción de datos}

Para cada estudio incluido se extrajeron autor, año, país, objetivo, población, tipo de datos, instrumentos, variable principal, método de análisis, resultados, limitaciones y aporte a la revisión. La síntesis fue narrativa y comparativa, debido a la heterogeneidad de métodos, escalas e indicadores.


\subsection{Evaluación y síntesis de evidencia}

La evaluación de la evidencia se realizó mediante una síntesis narrativa y comparativa. No se aplicó metaanálisis debido a la heterogeneidad de los estudios incluidos, ya que estos utilizaron distintas poblaciones, instrumentos, diseños metodológicos, variables principales y métricas de evaluación.

Los artículos fueron comparados según tipo de población, variable asociada al burnout académico, tipo de datos, método de análisis, principales resultados y limitaciones. Para organizar la evidencia se agruparon los estudios por tipo de método, tipo de dato, factores de riesgo, factores protectores y vacíos metodológicos.

También se distinguió entre evidencia directa y complementaria. La evidencia directa correspondió a estudios centrados en burnout académico, agotamiento o cansancio emocional en estudiantes universitarios. La evidencia complementaria incluyó trabajos sobre salud mental estudiantil, hábitos de estudio, sueño, uso de tecnología, estrés institucional o analítica educativa, siempre que aportaran variables o métodos útiles para el modelado del riesgo académico.

ChatGPT fue utilizado como apoyo para ordenar la síntesis, identificar categorías comparativas y mejorar la redacción académica. Sin embargo, la interpretación final fue revisada manualmente con base en los artículos incluidos.

\subsection{Uso de herramientas computacionales}

Durante la revisión se utilizaron herramientas computacionales para apoyar la limpieza, organización, análisis y presentación de los registros bibliográficos.

Python fue utilizado para integrar los archivos exportados desde Scopus y Web of Science, normalizar campos, detectar duplicados, aplicar filtros de cribado, generar archivos CSV/XLSX y construir el diagrama PRISMA. Las planillas procesadas permitieron revisar registros, documentar decisiones de inclusión o exclusión y construir la matriz de extracción de datos.

LaTeX fue utilizado para redactar y compilar el manuscrito final en PDF. Además, se utilizó ChatGPT como apoyo metodológico y de escritura para estructurar secciones, depurar criterios, organizar la síntesis de evidencia, revisar coherencia textual y formular borradores preliminares.

El uso de ChatGPT fue auxiliar y no reemplazó la revisión humana. La selección final de artículos, la recuperación de textos completos, la interpretación de resultados y la validación de la información fueron realizadas manualmente.

\subsection{Uso declarado de inteligencia artificial en la metodología}

Antes de presentar la tabla de uso de inteligencia artificial, es importante declarar cómo se integraron estas herramientas en el proceso metodológico. En esta revisión se utilizaron dos herramientas: Gemini y ChatGPT. Gemini fue empleado únicamente en la etapa inicial de formulación del problema, con el propósito de explorar posibles enfoques de investigación vinculados con salud mental, educación superior, burnout académico y Ciencia de Datos. Posteriormente, ChatGPT fue utilizado como apoyo auxiliar en la organización metodológica, extracción de datos, síntesis de evidencia, redacción académica y construcción de tablas.

Las instrucciones entregadas a las herramientas de IA fueron de carácter general y metodológico. Por ejemplo, se solicitaron apoyos para ordenar criterios de inclusión y exclusión, estructurar el flujo PRISMA, resumir hallazgos de los estudios incluidos, comparar métodos de análisis y mejorar la coherencia del manuscrito. No se utilizó IA para reemplazar la lectura crítica de los artículos ni para decidir automáticamente qué estudios debían ser incluidos o excluidos.

Las decisiones centrales fueron tomadas manualmente por el investigador. Esto incluye la definición final del tema, la selección de bases de datos, la aplicación de criterios de inclusión y exclusión, la recuperación de textos completos, la elección de los 6 estudios incluidos y la interpretación final de los resultados. La IA funcionó como una herramienta de apoyo para organizar y redactar, pero no como fuente primaria de evidencia.

La información generada con apoyo de IA fue verificada mediante contraste con los documentos originales, la base filtrada, el diagrama PRISMA, la matriz de extracción y los artículos seleccionados. En los casos donde la IA propuso formulaciones generales o interpretaciones amplias, estas fueron corregidas para mantener coherencia con la evidencia revisada.

Finalmente, se reconoce que el uso de IA presenta riesgos metodológicos, como simplificación excesiva, errores de interpretación, omisión de matices o generación de afirmaciones no respaldadas. Para reducir estos riesgos, todas las salidas generadas por IA fueron revisadas manualmente antes de incorporarse al manuscrito final.

\begin{table}[htbp]
\centering
\caption{Uso de inteligencia artificial y verificación humana}
\label{tab:uso_ia}
\renewcommand{\arraystretch}{1.15}
\small
\begin{tabularx}{\textwidth}{|
>{\raggedright\arraybackslash}p{2.6cm}|
>{\raggedright\arraybackslash}p{2.3cm}|
>{\raggedright\arraybackslash}X|
>{\raggedright\arraybackslash}X|}
\hline
\textbf{Etapa} & \textbf{Herramienta} & \textbf{Uso específico} & \textbf{Verificación humana} \\
\hline

Formulación del problema &
Gemini &
Exploración inicial del tema y posibles enfoques de investigación. &
Se ajustó el tema según la asignatura y la literatura disponible. \\
\hline

Búsqueda y criterios &
ChatGPT &
Apoyo para ordenar palabras clave, operadores booleanos y criterios de inclusión/exclusión. &
Las búsquedas y criterios fueron revisados y aplicados manualmente. \\
\hline

Extracción de datos &
ChatGPT &
Apoyo para estructurar la matriz de extracción: autor, año, población, método, variables, resultados y limitaciones. &
Los datos fueron contrastados manualmente con los textos completos de los artículos incluidos. \\
\hline

Síntesis de evidencia &
ChatGPT &
Agrupación de estudios por método, tipo de dato, hallazgos y vacíos. &
La síntesis fue contrastada con los 6 artículos incluidos. \\
\hline

Redacción y tablas &
ChatGPT &
Apoyo en redacción, coherencia, tablas y material suplementario. &
El texto y las tablas fueron revisados y corregidos manualmente. \\
\hline

\end{tabularx}
\end{table}

\section{Resultados}

\subsection{Proceso de selección}

La búsqueda identificó 1.020 registros: 886 en Web of Science y 134 en Scopus. Tras eliminar 2 duplicados, se cribaron 1.018 registros. En esta etapa se excluyeron 1.004 por no cumplir criterios de pertinencia. Posteriormente, se buscaron 14 informes para texto completo; 2 no fueron recuperados. De los 12 informes evaluados, 6 fueron excluidos con razón y 6 estudios fueron incluidos en la revisión.

\subsection{Estudios incluidos}

Los 6 estudios incluidos abordan burnout académico, agotamiento estudiantil, salud mental o factores asociados en contextos educativos (Alluri et al., 2026; Bojorque et al., 2025; Cuevas Caravaca et al., 2025; Mud\l{}o-G\l{}agolska \& Larionow, 2025; Rodrigues Matos et al., 2024; Rodríguez Muñoz \& Antino, 2021). Aunque existe heterogeneidad metodológica, todos aportan evidencia útil para analizar variables predictoras, factores protectores o técnicas aplicables desde Ciencia de Datos.

\begin{table}[H]
\centering
\caption{Estudios incluidos en la revisión sistemática}
\label{tab:estudios_incluidos}
\renewcommand{\arraystretch}{1.15}
\scriptsize

\begin{tabularx}{\textwidth}{|
>{\centering\arraybackslash}p{0.04\textwidth}|
>{\raggedright\arraybackslash}p{0.17\textwidth}|
>{\raggedright\arraybackslash}p{0.22\textwidth}|
>{\raggedright\arraybackslash}p{0.25\textwidth}|
>{\raggedright\arraybackslash}X|}
\hline

\textbf{N°} &
\textbf{Autor/año} &
\textbf{Población / datos} &
\textbf{Método principal} &
\textbf{Aporte a la revisión} \\
\hline

1 &
Alluri et al. (2026) &
Estudiantes de educación superior del Healthy Minds Network. &
K-means estándar, K-means justo, bootstrap, pruebas estadísticas y double machine learning. &
Segmentación de perfiles de salud mental con equidad algorítmica. \\
\hline

2 &
Cuevas Caravaca et al. (2025) &
789 universitarios de Educación Primaria y Enfermería. &
Correlación, t de Student y regresión lineal múltiple. &
Identifica el ejercicio físico como posible factor protector entre burnout y cansancio emocional. \\
\hline

3 &
Mudło-Głagolska y Larionow (2025) &
Tres muestras de estudiantes universitarios de primer año. &
Correlaciones, SEM longitudinal y análisis de red bayesiano. &
Muestra el rol protector de la pasión armoniosa frente a agotamiento y desenganche. \\
\hline

4 &
Bojorque et al. (2025) &
Ocho años de datos institucionales de una universidad ecuatoriana. &
Minería de datos y análisis exploratorio. &
Evidencia el uso de analítica institucional para detectar estrés y anomalías académicas. \\
\hline

5 &
Rodrigues Matos et al. (2024) &
401 estudiantes en transición hacia educación superior. &
Estadística descriptiva y análisis de redes. &
Relaciona exaustión, sueño, tiempo de estudio, ansiedad y desenganche. \\
\hline

6 &
Rodríguez Muñoz y Antino (2021) &
45 universitarios y 450 observaciones diarias. &
Diseño de diario y análisis multinivel. &
El uso del smartphone en clase predice menor engagement y mayor agotamiento diario. \\
\hline

\end{tabularx}
\end{table}

\subsection{Variables y métodos identificados}

\begin{table}[H]
\centering
\caption{Síntesis de variables y métodos identificados}
\label{tab:variables_metodos}
\renewcommand{\arraystretch}{1.15}
\scriptsize

\begin{tabularx}{\textwidth}{|
>{\raggedright\arraybackslash}p{0.17\textwidth}|
>{\raggedright\arraybackslash}p{0.27\textwidth}|
>{\raggedright\arraybackslash}p{0.22\textwidth}|
>{\raggedright\arraybackslash}X|}
\hline

\textbf{Dimensión} &
\textbf{Variables principales} &
\textbf{Métodos utilizados} &
\textbf{Interpretación} \\
\hline

Burnout y agotamiento &
Agotamiento, cinismo, eficacia, cansancio emocional, desenganche. &
Regresión, correlación, redes, multinivel. &
Núcleo del fenómeno; el agotamiento aparece como variable central. \\
\hline

Salud mental &
Ansiedad, depresión, bienestar, PHQ-9, GAD-7, flourishing. &
Redes, SEM, clustering, pruebas estadísticas. &
El burnout forma parte de un sistema amplio de salud mental. \\
\hline

Motivación académica &
Pasión armoniosa, pasión obsesiva, engagement, abandono. &
SEM, correlaciones, multinivel. &
La pasión armoniosa y el engagement operan como factores protectores. \\
\hline

Hábitos diarios &
Sueño, horas de estudio, smartphone, ejercicio físico. &
Redes, diario, regresión, multinivel. &
Variables modificables útiles para alerta temprana. \\
\hline

Variables sociodemográficas &
Sexo, edad, carrera, raza, nivel socioeconómico. &
Pruebas de medias, clustering justo, bootstrap. &
Pueden explicar riesgo, pero requieren tratamiento ético y equidad algorítmica. \\
\hline

Datos institucionales &
Matrícula, retención, notas, carga docente, empleabilidad. &
Minería de datos, EDA, visualización temporal. &
Permiten monitorear patrones agregados de estrés y rendimiento. \\
\hline

\end{tabularx}
\end{table}

\subsection{Síntesis de hallazgos}

Los estudios revisados muestran que el burnout académico se asocia con variables psicológicas, conductuales e institucionales. Entre los factores de riesgo destacan agotamiento, cinismo, desenganche, ansiedad, síntomas depresivos, menor sueño, mayor tiempo de estudio, uso del smartphone en clase y experiencias de discriminación (Alluri et al., 2026; Rodrigues Matos et al., 2024; Rodríguez Muñoz \& Antino, 2021). Entre los factores protectores aparecen ejercicio físico, pasión armoniosa, estabilidad financiera, sueño suficiente y participación social o cultural (Cuevas Caravaca et al., 2025; Mud\l{}o-G\l{}agolska \& Larionow, 2025).

Los métodos identificados incluyen regresión lineal, modelos multinivel, SEM, análisis de redes, minería de datos y clustering. No se observa predominio de modelos predictivos avanzados como Random Forest, XGBoost o redes neuronales. En cambio, predominan enfoques explicativos e interpretables. El estudio de clustering justo representa el enfoque computacional más avanzado, al incorporar segmentación de perfiles, equidad algorítmica e inferencia causal.

La comparación entre estudios se ve limitada por la diversidad de métricas. Las regresiones reportan coeficientes y varianza explicada; los modelos longitudinales usan índices de ajuste como CFI, TLI, RMSEA y SRMR; los modelos multinivel analizan variabilidad diaria; las redes usan centralidad o fuerza de conexión; y el clustering evalúa costos de agrupamiento y disparidades entre grupos. Por ello, la síntesis se realizó de forma narrativa y no mediante metaanálisis.


\subsection{Operacionalización para un modelo predictivo}

A partir de los estudios incluidos, es posible proponer una operacionalización preliminar del problema para futuras investigaciones. La variable objetivo podría definirse como riesgo de burnout académico alto, nivel de agotamiento académico, cansancio emocional elevado o pertenencia a un perfil de riesgo. La elección depende del instrumento disponible y del objetivo institucional. Si se usa una escala como MBI-SS u OLBI, la salida del modelo puede formularse como clasificación binaria, clasificación ordinal o predicción de puntaje continuo.

Las variables predictoras deberían organizarse en cuatro grupos. El primero corresponde a indicadores psicológicos, como ansiedad, depresión, bienestar, engagement, pasión por estudiar y desenganche. El segundo incluye variables académicas, como carrera, avance curricular, notas, asistencia, carga de asignaturas, reprobaciones o intención de abandono. El tercero agrupa hábitos y comportamiento diario, como horas de sueño, horas de estudio, ejercicio físico, uso de smartphone en clase y participación social. El cuarto considera variables sociodemográficas y contextuales, como edad, sexo, condición socioeconómica y experiencias de discriminación.

\begin{table}[H]
\centering
\caption{Propuesta de operacionalización para futuros modelos predictivos}
\label{tab:operacionalizacion_modelo}
\renewcommand{\arraystretch}{1.15}
\scriptsize

\begin{tabularx}{\textwidth}{|
>{\raggedright\arraybackslash}p{0.20\textwidth}|
>{\raggedright\arraybackslash}p{0.31\textwidth}|
>{\raggedright\arraybackslash}p{0.20\textwidth}|
>{\raggedright\arraybackslash}X|}
\hline

\textbf{Componente} &
\textbf{Ejemplos de variables} &
\textbf{Tipo de dato} &
\textbf{Uso potencial} \\
\hline

Variable objetivo &
Burnout alto, agotamiento, cansancio emocional, perfil \textit{at-risk}. &
Escala o clase. &
Definir la salida del modelo. \\
\hline

Psicológico &
Ansiedad, depresión, bienestar, engagement, pasión armoniosa. &
Encuesta validada. &
Estimar riesgo individual. \\
\hline

Académico &
Notas, avance curricular, carga académica, asistencia, reprobaciones. &
Registro institucional. &
Complementar evidencia subjetiva. \\
\hline

Conductual &
Sueño, horas de estudio, ejercicio, smartphone, participación social. &
Encuesta o registro diario. &
Detectar factores modificables. \\
\hline

Sociodemográfico &
Edad, sexo, carrera, nivel socioeconómico, grupo subrepresentado. &
Registro o encuesta. &
Evaluar equidad y contexto. \\
\hline

Evaluación &
AUC, F1-score, sensibilidad, especificidad, calibración, explicabilidad. &
Métrica de modelo. &
Comparar desempeño y utilidad. \\
\hline

\end{tabularx}
\end{table}

Esta propuesta muestra que el problema no debe reducirse a un algoritmo específico. Un sistema útil debería combinar desempeño predictivo con interpretabilidad y criterios de equidad. Además, la evaluación no debería limitarse a accuracy, especialmente si existe desbalance entre estudiantes con bajo y alto riesgo. Métricas como sensibilidad, especificidad, F1-score, AUC, calibración y análisis de falsos positivos resultan más adecuadas para un contexto preventivo.

\section{Discusión}

Los estudios revisados muestran que el burnout académico no es un fenómeno aislado, sino un sistema de variables psicológicas, conductuales, académicas e institucionales. El agotamiento se vincula con cansancio emocional, ansiedad, desenganche, menor engagement, bajo bienestar y mayor vulnerabilidad frente a factores de estrés. En contraste, el ejercicio físico, la pasión armoniosa, el sueño suficiente y la participación social aparecen como factores protectores.

Desde el punto de vista metodológico, los trabajos se agrupan en tres enfoques. El primero corresponde a modelos explicativos tradicionales, como correlaciones y regresiones, útiles por su interpretabilidad, pero limitados para capturar relaciones complejas. El segundo incluye diseños longitudinales, multinivel y de diario, que permiten observar cambios temporales e intraindividuales, como el efecto diario del uso del smartphone sobre agotamiento. El tercero incorpora métodos más cercanos a Ciencia de Datos, como clustering, análisis de redes y minería de datos institucional, capaces de representar perfiles, relaciones multivariadas y patrones agregados.

La revisión muestra que no existe una técnica dominante. La regresión permite identificar predictores interpretables; los modelos multinivel capturan variación diaria; el SEM evalúa relaciones temporales; las redes identifican variables centrales; el clustering segmenta perfiles de riesgo; y la minería de datos institucional permite observar patrones en sistemas educativos completos. Para una solución aplicada, lo más adecuado sería integrar varias técnicas en vez de depender de un único modelo.

Un hallazgo importante es la escasa presencia de modelos predictivos avanzados y reproducibles. Aunque los estudios aplican métodos analíticos relevantes, predominan enfoques explicativos e inferenciales. Esto deja una oportunidad concreta para Ciencia de Datos: desarrollar modelos predictivos robustos, interpretables y validados externamente para apoyar sistemas de alerta temprana en educación superior.

También se identifican vacíos relevantes. Existe baja disponibilidad de datasets abiertos y código; las muestras suelen ser locales; los instrumentos de medición del burnout no son homogéneos; y hay poca integración entre datos psicológicos, académicos, conductuales e institucionales. Además, pocos estudios incorporan explícitamente criterios de equidad algorítmica, pese a que variables como género, carrera, raza o nivel socioeconómico pueden influir en la clasificación de riesgo.

Una futura línea de investigación debería diseñar modelos integrados de detección temprana que combinen encuestas validadas, registros académicos, hábitos de sueño, carga de estudio, uso de tecnología, engagement y variables sociodemográficas. Estos modelos deben ser explicables, auditables y usados con fines preventivos, no punitivos. En problemas de bienestar estudiantil, la precisión predictiva no basta: también se requiere privacidad, consentimiento informado, gobernanza de datos y revisión humana.

En síntesis, la literatura muestra potencial para aplicar Ciencia de Datos al burnout académico, pero el campo aún está en consolidación. Hay evidencia sobre factores asociados y métodos útiles, pero falta avanzar hacia modelos predictivos transferibles, reproducibles, éticos y aplicables en contextos reales de educación superior.

\section{Limitaciones}

Esta revisión presenta varias limitaciones. Primero, la búsqueda se realizó solo en Web of Science y Scopus, por lo que pudieron quedar fuera estudios indexados en bases como PubMed, ERIC, PsycINFO, IEEE Xplore, ACM o Google Scholar.

Segundo, la estrategia dependió de términos específicos sobre burnout académico, estudiantes universitarios y Ciencia de Datos. Esto pudo excluir estudios que usan conceptos cercanos, como estrés académico, bienestar, fatiga, engagement o intención de abandono, aunque también permitió reducir ruido temático.

Tercero, 2 de los 14 informes candidatos no pudieron recuperarse a texto completo, lo que introduce sesgo de disponibilidad. Además, de los 12 informes evaluados, 6 fueron excluidos por no responder suficientemente a la pregunta de investigación.

Cuarto, los estudios incluidos son metodológicamente heterogéneos. Usan regresión, análisis multinivel, SEM, redes, minería de datos y clustering, lo que impide comparar resultados con una métrica común o realizar metaanálisis.

Quinto, no todos los estudios miden burnout académico con el mismo instrumento. Algunos usan MBI-SS, otros OLBI adaptado, y otros trabajan con cansancio emocional, agotamiento o desenganche. Esto obliga a interpretar los resultados como evidencia relacionada, no como mediciones completamente equivalentes.

Sexto, algunos estudios aportan evidencia indirecta, especialmente aquellos centrados en estrés institucional, hábitos de estudio, sueño o salud mental general. Fueron incluidos porque contribuyen al análisis desde Ciencia de Datos, pero no todos tienen el mismo nivel de alineación con la pregunta principal.

Finalmente, la revisión utilizó herramientas computacionales e IA como apoyo. Aunque la selección e interpretación fueron revisadas manualmente, el uso de IA puede introducir simplificaciones o errores si no se verifica críticamente. Tampoco se aplicó una escala formal de riesgo de sesgo, lo que debe considerarse en futuras revisiones.

\section{Conclusiones}

Esta revisión sistemática analizó el uso de técnicas de Ciencia de Datos para la predicción o modelado del burnout académico en estudiantes universitarios. A partir de 1.020 registros identificados en Web of Science y Scopus, el proceso PRISMA permitió incluir finalmente 6 estudios.

Los métodos encontrados incluyen regresión lineal múltiple, análisis multinivel, modelos de ecuaciones estructurales, análisis de redes, minería de datos institucional y clustering. Los enfoques más cercanos a Ciencia de Datos corresponden al clustering, las redes y la minería de datos, ya que permiten analizar múltiples variables, construir perfiles de riesgo y detectar patrones en sistemas educativos.

La evidencia indica que el burnout académico se relaciona con factores psicológicos, conductuales, académicos, sociodemográficos e institucionales. Variables como agotamiento, cinismo, desenganche, ansiedad, menor sueño, mayor tiempo de estudio y uso del smartphone en clase aparecen asociadas a mayor riesgo. En cambio, la práctica de ejercicio físico, la pasión armoniosa, el sueño suficiente y la participación social aparecen como posibles factores protectores.

El principal vacío identificado es la falta de modelos predictivos avanzados, reproducibles y validados externamente. Predominan enfoques explicativos, mientras que aún hay escasa evidencia sobre modelos aplicables a sistemas reales de alerta temprana. También se observa baja disponibilidad de datos abiertos, poca integración entre registros psicológicos e institucionales, y limitada consideración de equidad algorítmica.

Como proyección, futuras investigaciones deberían construir modelos integrados que combinen encuestas validadas, datos académicos, hábitos de estudio, sueño, uso de tecnología y variables institucionales. Estos modelos deben ser interpretables, auditables y éticamente responsables, con consentimiento informado, protección de datos y uso preventivo.

En conclusión, la Ciencia de Datos tiene alto potencial para apoyar la detección temprana y prevención del burnout académico, pero el campo requiere mayor madurez metodológica, reproducibilidad, validación externa y conexión con la toma de decisiones institucional.


\section*{Referencias}
\addcontentsline{toc}{section}{Referencias}

\apaRef{Alluri, P., Chen, Z., Thesen, T., Jacobson, N. C., \& Marrero, W. J. (2026). Fairness-aware K-means clustering in digital mental health for higher education students: A generalizable framework for equitable clustering. \textit{JAMIA Open, 9}(1), ooaf174. \url{https://doi.org/10.1093/jamiaopen/ooaf174}}

\apaRef{Bojorque, R., Moscoso, F., Pesántez, F., \& Flores, Á. (2025). Stress factors in higher education: A data analysis case. \textit{Data, 10}, 22. \url{https://doi.org/10.3390/data10020022}}

\apaRef{Cuevas Caravaca, E., Garcés de los Fayos Ruiz, E. J., Sánchez Romero, E. I., \& López Morales, J. L. (2025). Academic burnout and emotional exhaustion among university students: The role of physical exercise. \textit{SPORT TK-EuroAmerican Journal of Sport Sciences, 14}, Article 99.}

\apaRef{Mud\l{}o-G\l{}agolska, K., \& Larionow, P. (2025). Passion for studying and its relationships with academic burnout and mental health: Longitudinal insights into sustainable students' functioning. \textit{Sustainability, 17}, 9852. \url{https://doi.org/10.3390/su17219852}}

\apaRef{Page, M. J., McKenzie, J. E., Bossuyt, P. M., Boutron, I., Hoffmann, T. C., Mulrow, C. D., Shamseer, L., Tetzlaff, J. M., Akl, E. A., Brennan, S. E., Chou, R., Glanville, J., Grimshaw, J. M., Hróbjartsson, A., Lalu, M. M., Li, T., Loder, E. W., Mayo-Wilson, E., McDonald, S., ... Moher, D. (2021). The PRISMA 2020 statement: An updated guideline for reporting systematic reviews. \textit{BMJ, 372}, n71. \url{https://doi.org/10.1136/bmj.n71}}

\apaRef{Rodrigues Matos, F., Malavasi, R. M., \& De Andrade, A. L. (2024). Relações entre burnout estudantil, saúde mental, hábitos de estudo e sono: Um estudo transversal. \textit{Pensando Psicología, 20}(1), 1--19. \url{https://doi.org/10.16925/2382-3984.2024.01.01}}

\apaRef{Rodríguez Muñoz, A., \& Antino, M. (2021). El uso del teléfono móvil en clase y su efecto sobre el engagement académico y el agotamiento: Un estudio de diario en estudiantes universitarios. \textit{European Journal of Education and Psychology, 14}(1), 1--10. \url{https://doi.org/10.32457/ejep.v14i1.1401}}

\end{document}
